%!TEX root = ../main.tex

\addchap{Anhang}
\begin{itemize}
	\item \ref{apx:anforderungen} Anforderungen \dotfill{} \pageref{apx:anforderungen}
	\item \ref{apx:konzept_architektur} Konzeptionierte Architektur \dotfill{} \pageref{apx:konzept_architektur}
	\item \ref{apx:implementierteseiten} Implementierte Seiten \dotfill{} \pageref{apx:implementierteseiten}
	\item \ref{apx:optischetests} Optische Tests \dotfill{} \pageref{apx:optischetests}
	\item \ref{apx:softwaretests} Software Tests \dotfill{} \pageref{apx:softwaretests}
\end{itemize}

\chapter{Anforderungen}
\label{apx:anforderungen}
\begin{figure}[H]
    \includepdf[width=0.9\textwidth]{Anforderungen Moscow.pdf}
\end{figure}

\chapter{Konzeptionierte Architektur}
\label{apx:konzept_architektur}
\begin{figure}[H]
    \centering
    \includegraphics[width=\textwidth]{ARis_Systemarchitektur.jpg}
    \caption{Software-Architektur der AR-Brille}
\end{figure}

\chapter{Implementierte Seiten}
\label{apx:implementierteseiten}
\begin{figure}[H]
    \centering
    \includegraphics[width=\textwidth]{HUD.jpeg}
    \caption{Hauptseite mit verschiedenen Modulen}
\end{figure}
\begin{figure}[H]
    \centering
    \includegraphics[width=\textwidth]{Teleprompter.jpeg}
    \caption{Startseite des Teleprompters, wartet auf das Signal der Companion App}
\end{figure}
\begin{figure}[H]
    \centering
    \includegraphics[width=\textwidth]{Navigation1.jpeg}
    \caption{Mögliche Navigationsansicht}
\end{figure}
\begin{figure}[H]
    \centering
    \includegraphics[width=\textwidth]{Navigation2.jpeg}
    \caption{Mögliche Navigationsansicht}
\end{figure}
\begin{figure}[H]
    \centering
    \includegraphics[width=\textwidth]{Navigation3.jpeg}
    \caption{Mögliche Navigationsansicht}
\end{figure}

\chapter{Optische Tests}
\label{apx:optischetests}
\begin{table}[htbp]
  \centering
  \caption{Ergebnisse der optischen Komponententests.}
  \begin{tabularx}{\textwidth}{@{} l X l @{}}
    \toprule
    \textbf{Prüfpunkt} & \textbf{Methode und Befund} & \textbf{Ergebnis} \\
    \midrule
    Display-Funktion &
      Spannungsversorgung angelegt; Display leuchtet gleichmäßig auf.
      Funktion im weiteren Betrieb bestätigt. &
      Erfolgreich \\
    \addlinespace
    Display-Linsen-Abstand &
      Ausgangsabstand des originalen Viewfinders als Startwert verwendet;
      iterative Anpassung bis zum schärfsten Bild.
      Scharfe Abbildung im eingestellten Abstand erreicht. &
      Erfolgreich \\
    \addlinespace
    Spiegelwinkel &
      Spiegel wurden auf ihre
      konstruktiv vorgegebenen Winkel überprüft und justiert.
      Korrekte Strahlführung zum Auge bestätigt. &
      Erfolgreich \\
    \addlinespace
    See-Through-Funktion &
      Durchblick auf reale Umgebung durch den halbdurchlässigen
      Combiner-Spiegel bei gleichzeitiger Display-Anzeige qualitativ
      geprüft. Überlagerung visuell erkennbar. &
      Erfolgreich \\
    \addlinespace
    Bildgröße / FoV &
      Das auf dem Combiner sichtbare Bild ist deutlich kleiner als
      konzeptionell angestrebt. Die geringe Vergrößerung des
      Viewfinder-Okulars sowie der kurze optische Weg zwischen
      Display und Combiner führen dazu, dass nur ein
      Ausschnitt des Displays im Spiegel erscheint
      (vgl.~Abschnitt~\ref{sec:optischeImplementierung}). &
      Teilweise erfolgreich \\
    \addlinespace
    Fokussierung &
      Die Schärfentiefe des Systems erlaubt keine simultane Fokussierung
      auf Display-Inhalt und reale Umgebung. Der Nutzer muss aktiv
      zwischen Nah- und Fernfokus wechseln, was dem Ziel einer
      entspannten Augenakkommodation widerspricht. &
      Teilweise erfolgreich \\
    \bottomrule
  \end{tabularx}
\end{table}

\chapter{Software Tests}
\label{apx:softwaretests}
\begin{table}[htbp]
  \centering
  \caption{Ergebnisse der Software-Komponenten- und Systemtests.}
  \begin{tabularx}{\textwidth}{@{} l X l @{}}
    \toprule
    \textbf{Testbereich} & \textbf{Befund} & \textbf{Ergebnis} \\
    \midrule
    Backend-Start \& API &
      Flask-Backend startet stabil auf dem Raspberry~Pi; alle
      REST-Endpunkte (vgl.~Tabelle~\ref{tab:endpoints}) antworten
      korrekt auf HTTP-Anfragen vom lokalen Gerät. &
      Erfolgreich \\
    \addlinespace
    Sensorik (MPU-6050, GY-63) &
      Beide I2C-Sensoren werden beim Start korrekt initialisiert;
      Messwerte für Temperatur, Luftfeuchtigkeit und
      Beschleunigung werden über \texttt{/api/sensors} abrufbar. &
      Erfolgreich \\
    \addlinespace
    Taster-Eingabe &
      Single-Tap und Double-Tap werden von der Finite-State-Machine
      zuverlässig unterschieden; Seitenwechsel im Hub-Kontext
      funktioniert wie erwartet. &
      Erfolgreich \\
    \addlinespace
    SvelteKit-Frontend / Hub &
      Frontend startet als lokaler Webserver; Hub-Seite wird
      korrekt gerendert und zeigt Header. &
      Erfolgreich \\
    \addlinespace
    Demo-Seite &
      Kompakte Demo-Seite zeigt Kerninformationen auf minimalem
      Anzeigebereich; Taster-Interaktion funktioniert.
      Erreichtes Ziel: Proof-of-Concept-Anzeige auf der physischen
      Hardware. &
      Erfolgreich \\
    \addlinespace
    Teleprompter &
      Konfigurationsübertragung von der Companion App zum Backend
      und Polling durch das Frontend funktionieren;
      scrollende Textdarstellung im Browser verifiziert.
      Vollständiger End-to-End-Test auf der finalen Hardware
      nicht abgeschlossen. &
      Teilweise erfolgreich \\
    \addlinespace
    KI-Chat &
      Nachrichten können über die REST-API gespeichert und auf
      der Messages-Seite angezeigt werden. Weitere Implementation ist nicht abgeschlossen &
      Teilweise Erfolgreich \\
    \addlinespace
    Navigation &
      Seite strukturell vorhanden; MPU-6050-Daten werden über
      den Backend-Endpunkt bereitgestellt.
      Eine vollständige Navigationsfunktion (Turn-by-Turn o.\,ä.)
      wurde nicht implementiert. &
      Teilweise erfolgreich \\
    \addlinespace
    Companion App (WiFi) &
      REST-Kommunikation zwischen React-Native-App und Backend
      über WiFi funktioniert. Weitere Funktionalitäten wurden nicht getestet &
      Teilweise Erfolgreich \\
    \addlinespace
    Bluetooth / GATT-Server &
      GATT-Server-Adapter ist in der Onboard-Runtime implementiert;
      vollständiger BLE-Verbindungstest mit Companion App nicht
      abgeschlossen. Jedoch wurde Prototypische Verbindung zu einem Companion gerät hergestellt &
      Teilweise erfolgreich \\
    \addlinespace
    Pixel Streaming &
      Endpunkt im Backend angelegt; keine Implementierung oder
      Test durchgeführt. &
      Nicht umgesetzt \\
    \addlinespace
    Grenzfälle / Fehlerverhalten &
      Backend gibt bei fehlender Hardware (keine I2C-Sensoren)
      strukturierte Fehlerobjekte zurück, anstatt abzustürzen.
      Fehlende Companion-App-Verbindung wird toleriert. &
      Erfolgreich \\
    \bottomrule
  \end{tabularx}
\end{table}